\scanpage{146}
aux hom
\[
\bar d^0:P^0\longrightarrow K^1
\]
satisfaisant
\[
\bar d^0(xf)=\bar d^0(x)d^0(f)\qquad x\in P^0(\cdot),\ f\in K^0(\cdot).
\]
La condition mise ci-dessus sur $\zeta$ signifie simplement que $\bar d^0$ prend ses valeurs dans $Z^1K=\operatorname{Ker}(d^2)$.

Ceci s'applique notamment à la description
de $H^1_{\mathrm{DR}}(S,E)$, comme classifiant les objets suivants.

Un torseur $P$ sous $E$, un hom de faisceaux
d'ensembles
\[
d_P^{\deleted{\unsure{0}}}:P\longrightarrow\deleted{\unsure{Z^1(|Z(}}E\otimes\Omega^1_{S/T}\deleted{\unsure{)}}, \tag*{(17)}
\]
satisfaisant les conditions
\[
\left\{\begin{aligned}
d_P^{\deleted{\unsure{0}}}(xf)&=d_P^{\deleted{\unsure{0}}}(x)d_E^0(f)
 &x&\in P(\cdot),\ f\in\underline O_S(\cdot).\\
d_P(P)&\subset Z^1(E\otimes\Omega^{\unsure{1|\cdot|{}}}_{S/T})\ \text{i.e.}\ d_E^1d_P=0
\end{aligned}\right. \tag*{(18)}
\]
Un tel $d_P$ sur un $P$ s'appellera une \uline{connexion} \inserted{\uline{intégrable}} \uline{sur
le torseur} $P$, \uline{compatible avec la connexion} \inserted{\uline{intégrable}} $d_E^0$ \uline{sur} $E$.

\uline{NB} Si $d_E^0$ est une connexion sur $E$ (pas néc. intégrable),
\deleted{on peut \unsure{se pro}} on peut se proposer de trouver, pour un torseur $P$ sous $E$,
quelles sont les connexions compatibles avec $d_E^0$, au sens
général \deleted{\unsure{du ca}} de données de recollement infinitésimal
d'ordre $1$. Il est clair que si $S$ est affine, il existe
de telles connexions, donc il en existe tjrs localemt,
et que les connexions forment un torseur sous
$E\otimes\Omega^1$, qui en fait est isomorphe à $P\mathbin{\overset E{\times}}(E\otimes\Omega^1)$
(grâce à $d_E^0:E\longrightarrow E\otimes\Omega^1$). Donc les connexions
% archive sheet number: 145
